% Mục 6.3 (bản rút gọn). Rút từ sections/results_counter.tex; nguồn từng số ghi ở đó (results/apple-m5/counter.csv;
% scripts/report_numbers/rcross_cnt.py, rcross_cy.py, rbc_derived2.py, rbc_derived3.py).
\subsubsection{RQ5: Thứ tự hiệu năng của các counter}

\begin{rqanswer}\textbf{Trả lời RQ5.} Khi không có việc giữa hai lần tăng, fetch-and-add nhanh hơn CAS
2,0 tới 2,7 lần từ 4 thread, vì CAS hỏng và phải thử lại. Counter dùng spinlock không backoff chậm
nhất. Ở chế độ này, từ 4 thread, \PCname{lock_std_mutex} và \PCname{lock_pthread_mutex} ngang hoặc nhanh hơn
fetch-and-add. CAS có backoff nhanh hơn fetch-and-add từ 2 thread. CAS với sleep backoff cho các thread
tăng lần lượt (độ công bằng xấp xỉ $1/p$). \PCname{cas_spin} và \PCname{cas_yield} chỉ dồn một phần các
lần tăng về một thread (độ công bằng 0,52 và 0,39 ở 4 thread). Với mutex thư viện, cơ chế tương tự mới là
giả thuyết. Lợi thế này mất khi có việc giữa hai lần tăng. Số liệu của \PCname{apple-m5}. RQ6 kiểm lại
trên hai máy còn lại.\end{rqanswer}
% nguồn: bảng countertime apple-m5 (OW=0); thr_min/thr_max OW=0 p=4: cas_spin 0,52, cas_yield 0,39.

\PCfigure{apple-m5}{counter_time}

\emph{Bằng chứng.} Hình~\ref{fig:apple-m5:counter_time} vẽ thời gian chạy, quy về $10^6$ lần tăng mỗi
thread. Không counter nào mở rộng. Với $\mathrm{OW} = 0$, thông lượng của \PCname{faa} giảm từ khoảng 610
triệu lần tăng mỗi giây với 1 thread xuống 42 triệu với 20 thread, và từ 8 thread một lần tăng tốn gần
như không đổi (23--30~ns). \PCname{cas} chậm hơn \PCname{faa} 2,0 tới 2,7 lần từ 4 thread, nhưng ở 2
thread lại tốn ít hơn 28\%. Gọi $f$ là số lần CAS hỏng trên mỗi lần tăng
(Hình~\ref{fig:apple-m5:counter_cas_failures}). Ở $\mathrm{OW} = 0$, $f$ tăng đều từ 0,09 ở 2 thread lên
1,59 ở 8 thread. Ba biến thể CAS có backoff hỏng ít hơn \PCname{cas} một tới ba bậc ($f \le 0{,}02$), và ở
2 thread chúng chỉ tốn 3,4--3,9~ms cho $10^6$ lần tăng mỗi thread, so với 11,7~ms của \PCname{faa}.
% nguồn: faa OW=0, seconds/ops: p=8..20 23,4..30,2 ns; t(cas)/t(faa) OW=0 p=2 0,72; p=4..20 2,01..2,71;
% cas_fail/ops OW=0: p=2 0,091, p=8 1,59; backoff <= 0,022. Bảng countertime apple-m5 OW=0 p=2:
% cas_sleep/cas_spin/cas_yield 0,00341/0,00385/0,00340 s, faa 0,0117 s.

\PCfigure{apple-m5}{counter_cas_failures}

Ở $\mathrm{OW} = 0$ và 20 thread, \PCname{lock_tas} chậm hơn \PCname{faa} 13,7 lần, và
\PCname{lock_mcs} chậm hơn \PCname{lock_tas} 3,2 lần ngay ở 2 thread. Ngược lại, counter dùng
\PCname{lock_pthread_mutex} và \PCname{lock_std_mutex} ngang \PCname{faa} ở 4 thread (9,3 và 10,1 so với
10,2~ns) và nhanh hơn từ 6 thread. Ở 8 thread, một lần tăng tốn 12,4 và 14,1~ns, so với 23,4~ns của \PCname{faa}, và độ công bằng của hai mutex là
0,90--0,99. Ở 3 thread, chúng chậm hơn ở 4 thread (14,3--14,8 so với 9,3--10,1~ns), và nhóm chưa giải
thích được điều này. Ở $\mathrm{OW} = 1$, từ 8 thread mutex chậm hơn \PCname{faa} 3,0--3,4 lần. Với
$\mathrm{OW} = 8$, mọi counter gần nhau tới 4 thread, và \PCname{cas} chậm hơn \PCname{faa} 1,2--1,5 lần
từ 10 thread.
% nguồn: OW=0: lock_tas/faa p=20 13,67; lock_mcs/lock_tas p=2 3,23; ns/op p=3 std_mutex 14,3, pthread_mutex 14,8;
% p=4 10,05/9,3; p=8 pthread_mutex 12,41, std_mutex 14,10, faa 23,39; thr_min/thr_max 0,90-0,99; OW=1 p=8..20
% mutex/faa 3,00-3,42; OW=8 cas/faa p=10..20 1,21..1,45.

\emph{Vì sao.} Mỗi lần fetch-and-add phải giữ cache line của counter ở trạng thái độc quyền, nên các lần
tăng của mọi thread diễn ra lần lượt trên một line, và thêm thread không tăng thông lượng. Từ 8 thread,
một lần tăng tốn gần như không đổi, dù chưa chạm số core. Nhóm cho rằng line khi đó đã bận liên tục, nên
mỗi lần tăng tốn đúng một lần chuyển line, bất kể bao nhiêu thread đang chờ. Đây là giả thuyết. Một lần
CAS hỏng tốn một lần kéo cache line về mà không tăng được counter. Nếu mỗi lần hỏng tốn một lần chuyển, thời gian của \PCname{cas} xấp xỉ
thời gian của \PCname{faa} nhân $1+f$. Từ 6 tới 20 thread, mô hình lệch tỉ số đo được không quá 18\%. Ở 4
thread, mô hình chỉ giải thích 1,52 trong 2,14 lần.
% nguồn: (1+f) so với t(cas)/t(faa), p=4: 1,52/2,14; p=6..20 lệch <= 18 % (rbc_derived2.py).

Mô hình $1+f$ không giải thích được vì sao \PCname{cas} rẻ hơn \PCname{faa} ở 2 thread. Kiến trúc ARMv8.1
cho phép phần cứng thực hiện \texttt{ldadd} ở gần cache dùng chung thay vì trong L1 của core gọi. Nếu M5
làm vậy, mỗi lần \texttt{ldadd} tốn một lượt đi về tới cache dùng chung, còn một thread chạy vòng CAS có
thể tăng vài lần liên tiếp trên line đang nằm sẵn trong L1 của nó. Đây là giả thuyết. RQ6 cho thêm bằng
chứng gián tiếp từ hai máy khác.

\PCname{cas_sleep} cho các thread tăng lần lượt. Thời gian một lần tăng của nó bằng giá trị 1 thread
trong khoảng 11\% ở mọi số thread và mọi $\mathrm{OW}$. Vì vậy, ở 20 thread, \PCname{cas_sleep} nhanh hơn
\PCname{cas} 32 lần khi $\mathrm{OW} = 0$ nhưng chậm hơn 3,1 lần khi $\mathrm{OW} = 8$, vì chạy tuần tự
bỏ phí phần \texttt{is\_prime} lẽ ra chạy song song. Trong counter dùng lock, mỗi lần tăng chuyển ít nhất
hai line, lock và giá trị, và các thread chờ còn tranh line của lock, phù hợp với việc spinlock không
backoff chậm nhất. \PCname{lock_mcs} trao lock theo thứ tự đến, nên mỗi lần tăng chuyển cả lock lẫn giá
trị sang core khác. Với mutex thư viện, nhóm cho rằng người giữ lock thực hiện nhiều lần tăng liên tiếp
trên line đã nằm trong cache của mình, còn các thread chờ ngủ trong kernel và không tranh line. Nhóm chưa
đo độ dài các chuỗi này. Khi có việc giữa hai lần tăng, người nhả lock đi tính \texttt{is\_prime}, nên
lần lấy kế tiếp thường phải đánh thức một thread đang ngủ.
% nguồn: cas_sleep seconds/ops lệch p=1 tối đa 1,107; cas/cas_sleep OW=0 p=20 31,7, cas_sleep/cas OW=8 p=20 3,07.

Thứ tự \texttt{seq\_cst} làm \PCname{faa} dao động trong khoảng $\pm 23\%$, không theo hướng nào, nhưng
làm \PCname{cas} chậm khoảng 2 lần ở $\mathrm{OW} = 0$. Ở $\mathrm{OW} \ge 1$, khác biệt biến mất. Nhóm
giả thuyết rằng lệnh \texttt{casal} (acquire và release) ngăn các lần lặp chồng lên nhau.
% nguồn: t(cas_seqcst)/t(cas) OW=0: 1,29-2,61; OW=1,8: 0,83-1,10; t(faa_seqcst)/t(faa): 0,78-1,23.

\emph{Độ vững.} Nhóm không phân tích \PCname{lock_mcs} trên 10 thread, vì các ô đó đo với rất ít vòng
lặp. Chi phí \texttt{seq\_cst} của CAS không xuất hiện nhất quán trên \PCname{C4M7Y670NT} (RQ6), nên kết
luận về nó chỉ dựa trên \PCname{apple-m5}.
